%% guide-skieur-pente — la piste, le tire-fesses (montée) puis le départ de la descente.
%% Triangle rectangle : hypoténuse = piste de 50,0 m inclinée de 15°, dénivelé h à calculer.
%% RÈGLE : aucun vecteur poids ni réaction normale (étape 2), aucune valeur de h (étape 3),
%% aucune énergie ni travail chiffrés.
\documentclass[tikz,border=4pt]{standalone}
\usepackage{liaisons}
\begin{document}
\begin{tikzpicture}[x=1.15cm,y=1.15cm,
  sol/.style={line width=0.9pt, line join=round},
  skieur/.style={line width=1.1pt, draw=black!80, line cap=round},
  cote/.style={{Stealth[length=4pt]}-{Stealth[length=4pt]}, line width=0.7pt, draw=solideD},
  attache/.style={line width=0.35pt, draw=black!50},
  vect/.style={-{Stealth[length=6pt]}, line width=1.2pt},
  pet/.style={font=\scriptsize, inner sep=1.5pt}]

\def\al{15}      % inclinaison de la piste (degrés)

%% ================= la montagne ============================================
\filldraw[fill=black!6, draw=black!70, sol] (0,0) -- (6.0,0) -- (6.0,1.608) -- cycle;

%% angle à la base
\fill[solideD!25] (0,0) -- (0.95,0) arc (0:\al:0.95) -- cycle;
\draw[line width=0.5pt, draw=solideD] (0.95,0) arc (0:\al:0.95);
\node[pet, text=solideD] at (1.28,0.14) {15$^\circ$};

%% dénivelé h : côté vertical, en pointillé
\draw[cote, dash pattern=on 3pt off 2pt] (6.0,0.05) -- (6.0,1.558);
\node[font=\small, text=solideD, anchor=west, inner sep=3pt] at (6.06,0.80) {$h$};
\draw[line width=0.6pt, draw=black!70] (5.82,0) -- (5.82,0.18) -- (6.0,0.18);

%% ================= tout ce qui suit la pente ==============================
\begin{scope}[rotate=\al]
  %% câble du tire-fesses et ses deux poulies
  \draw[line width=0.8pt, draw=black!60] (0.35,1.15) -- (6.05,1.15);
  \draw[line width=0.7pt, fill=black!15] (0.35,1.15) circle (0.12);
  \draw[line width=0.7pt, fill=black!15] (6.05,1.15) circle (0.12);

  %% cote de la longueur de la piste
  \draw[attache] (0.10,0.02) -- (0.10,0.95);
  \draw[attache] (6.16,0.02) -- (6.16,0.95);
  \draw[cote] (0.10,0.85) -- (6.16,0.85);
  \node[pet, text=solideD, fill=white, inner sep=2pt] at (3.60,0.85) {\textbf{50,0 m}};

  %% skieur qui monte, accroché au câble
  \draw[skieur] (1.90,0) -- (1.90,0.42);
  \draw[skieur, fill=white] (1.90,0.55) circle (0.13);
  \draw[line width=0.5pt, draw=black!60] (1.90,0.68) -- (1.90,1.15);
  \draw[vect, draw=solideA] (2.15,0.28) -- (3.05,0.28);
  \node[font=\small, text=solideA, anchor=south, inner sep=1.5pt] at (2.60,0.32) {$\vec{T}$};

  %% skieur qui s'élance vers le bas de la piste
  \draw[skieur] (5.90,0) -- (5.90,0.42);
  \draw[skieur, fill=white] (5.90,0.55) circle (0.13);
  \draw[vect, draw=solideE] (5.55,0.28) -- (4.65,0.28);
  \node[pet, text=solideE, fill=white, inner sep=1.5pt] at (5.05,0.58) {$v_0 = 2{,}0$ m/s};
\end{scope}

%% ================= mentions ===============================================
\node[pet, text=black!70, anchor=north] at (1.90,-0.16) {montée à \textbf{vitesse constante}};
\end{tikzpicture}
\end{document}
